% !PRINTED 265
\weberpage{265}

\begin{center}
\textbf{\S\,74. The \gesperrt{Form} of the \gesperrt{Schlaefli} Modular Equations}\\
\textbf{for a Prime Degree.}
\end{center}

The form which we have found in the preceding section for the relations
subsisting between
\weq{1}{u=f(\omega),\qquad v=f\Bigl(\frac{c+\partial\omega}{a}\Bigr)}
% German source confirms the printed partial-shaped d glyph, transcribed as \partial.
can easily be brought under a general law, at least when the degree of
transformation is a prime $p$. Since much depends on this form in what
follows, we enter here somewhat more closely into it.

The determination of the numbers $r, s$ by (11) and (15) of the
preceding section depends on the behaviour of $p$ with respect to the
modulus 24, and since we may exclude the case $p=3$, we have the
following cases:
\weq{2}{\begin{array}{lll}
p\equiv 1 \ (\mathrm{mod}\ 24) & r=1 & s=12\\
p\equiv 5 & r=3 & s=2\\
p\equiv 7 & r=4 & s=3\\
p\equiv 11 & r=6 & s=1\\
p\equiv 13 & r=1 & s=6\\
p\equiv 17 & r=3 & s=4\\
p\equiv 19 & r=2 & s=3\\
p\equiv 23 & r=12 & s=1,
\end{array}}
so that $r+s$ is always odd, and $p+1$ is divisible by $2r$, $p-1$
by $2s$. Accordingly we obtain
\weq{3}{\begin{aligned}
A &= \Bigl(\frac{u}{v}\Bigr)^r+\Bigl(\frac{v}{u}\Bigr)^r,\\
B &= (uv)^s+\Bigl(\frac{2}{p}\Bigr)\frac{2^s}{(uv)^s}.
\end{aligned}}
Now we know that the $p+1$ quantities $v$ are the roots of an
irreducible transformation equation
\weq{4}{\Phi_p(v,u)=v^{p+1}+U_1v^p+\cdots U_{p+1}=0}
in which the coefficients $U_1\ldots, U_{p+1}$ are rational
functions of $u$.

We conclude at once that they are \gesperrt{polynomials} in $u$.
For in the \gesperrt{first place}, all the roots of (4) become
infinite for $u=\infty$, as one recognises on letting $\omega=i\infty$,
and therefore $q=0$. In the second place, by (9) of the preceding
%CONT>>
